% ==============================================================================
% Προοίμιο LaTeX (LaTeX Preamble)
% Πακέτα, Γραμματοσειρές και Ρυθμίσεις Μορφοποίησης
% ==============================================================================

% --- Γλώσσα και Γραμματοσειρές (XeLaTeX) ---
\usepackage{fontspec}
\usepackage{polyglossia}

\setmainlanguage{greek}
\setotherlanguage{english}

% Επιλογή γραμματοσειρών με πλήρη υποστήριξη Ελληνικών & Μαθηματικών
\setmainfont{Linux Libertine O}[
    Ligatures=TeX,
    Numbers=OldStyle
]
\setsansfont{Linux Biolinum O}[
    Ligatures=TeX
]
\setmonofont{DejaVu Sans Mono}[
    Scale=MatchLowercase
]

% --- Γεωμετρία Σελίδας ---
\usepackage[
    a4paper,
    top=2.8cm,
    bottom=2.8cm,
    left=3.0cm,
    right=2.5cm,
    headheight=15pt,
    footskip=1.2cm
]{geometry}

% Διάστιχο
\usepackage{setspace}
\setstretch{1.18}

% Ρύθμιση παραγράφων (παραδοσιακό πανεπιστημιακό στυλ)
\setlength{\parindent}{0.6cm}
\setlength{\parskip}{0.2em}

% Μείωση υπερχειλίσεων γραμμής (Overfull hbox) στο κύριο κείμενο
\tolerance=1500
\emergencystretch=1.5em
\hfuzz=2pt

% --- Μαθηματικά Πακέτα ---
\usepackage{amsmath}
\usepackage{amssymb}
\usepackage{amsfonts}
\usepackage{mathtools}
\usepackage{bm}

% --- Πίνακες & Σχήματα ---
\usepackage{graphicx}
\graphicspath{{figures/}{../evaluation/figures/}{../compose/}}
\usepackage{float}
\usepackage{booktabs}
\usepackage{tabularx}
\usepackage{multirow}
\usepackage{caption}
\usepackage{subcaption}
\usepackage{adjustbox}  % για κλιμάκωση πλατιών πινάκων

\captionsetup{
    font=small,
    labelfont=bf,
    labelsep=colon,
    figurename=Σχήμα,
    tablename=Πίνακας
}

% --- Χρώματα ---
\usepackage{xcolor}
\definecolor{ceidblue}{RGB}{0, 51, 102}
\definecolor{linkblue}{RGB}{20, 60, 160}
\definecolor{codebg}{RGB}{248, 248, 248}
\definecolor{codeframe}{RGB}{220, 220, 220}
\definecolor{commentgreen}{RGB}{40, 140, 50}
\definecolor{keywordblue}{RGB}{10, 40, 180}
\definecolor{stringred}{RGB}{180, 40, 30}

% --- Υπερσύνδεσμοι (Hyperref) ---
\usepackage{hyperref}
\hypersetup{
    colorlinks=true,
    linkcolor=linkblue,
    citecolor=linkblue,
    urlcolor=linkblue,
    pdftitle={\thesistitleel},
    pdfauthor={\thesisauthor},
    pdfsubject={\thesistype},
    bookmarksnumbered=true,
    bookmarksopen=true
}

% --- Κεφαλίδες & Υποσέλιδα (Fancyhdr) ---
\usepackage{fancyhdr}
\pagestyle{fancy}
\fancyhf{}
\fancyhead[LE,RO]{\thepage}
\fancyhead[RE]{\nouppercase{\leftmark}}
\fancyhead[LO]{\nouppercase{\rightmark}}
\renewcommand{\headrulewidth}{0.4pt}
\renewcommand{\footrulewidth}{0pt}

% Στυλ για κενές σελίδες στην αρχή κεφαλαίων
\fancypagestyle{plain}{
    \fancyhf{}
    \fancyfoot[C]{\thepage}
    \renewcommand{\headrulewidth}{0pt}
}

% Καθαρή λευκή σελίδα χωρίς κεφαλίδες όταν χρειάζεται
\newcommand{\clearemptydoublepage}{%
    \clearpage{\pagestyle{empty}\cleardoublepage}%
}

% --- Περιβάλλον Αλγορίθμων (Ruled Style όπως στο reference PDF) ---
\floatstyle{ruled}
\newfloat{algorithm}{tbp}{loa}[chapter]
\floatname{algorithm}{Αλγόριθμος}

% Βοηθητικές εντολές μορφοποίησης ψευδοκώδικα
\newcommand{\AlgoInput}[1]{\textbf{Input:} #1 \\}
\newcommand{\AlgoOutput}[1]{\textbf{Output:} #1 \\[0.15cm]}
\newcommand{\AlgoComment}[1]{\textit{/* #1 */} \\}
\newcommand{\AlgoLine}[2]{\textbf{#1} \quad #2 \\}
\newcommand{\AlgoLineIndent}[2]{\textbf{#1} \quad \quad #2 \\}
\newcommand{\AlgoLineIndentTwo}[2]{\textbf{#1} \quad \quad \quad #2 \\}

% --- Μορφοποίηση Πηγαίου Κώδικα (Listings) ---
\usepackage{listings}
\lstset{
    backgroundcolor=\color{codebg},
    basicstyle=\ttfamily\footnotesize,
    breaklines=true,
    captionpos=b,
    commentstyle=\color{commentgreen}\textit,
    frame=single,
    rulecolor=\color{codeframe},
    keywordstyle=\color{keywordblue}\textbf,
    stringstyle=\color{stringred},
    numbers=left,
    numberstyle=\tiny\color{gray},
    numbersep=8pt,
    tabsize=4,
    showstringspaces=false,
    keepspaces=true
}

% --- Βιβλιογραφία ---
\usepackage{cite}

% Προσαρμογή ελληνικών τίτλων
\addto\captionsgreek{
    \renewcommand{\contentsname}{Περιεχόμενα}
    \renewcommand{\listfigurename}{Κατάλογος Σχημάτων}
    \renewcommand{\listtablename}{Κατάλογος Πινάκων}
    \renewcommand{\bibname}{Βιβλιογραφία}
    \renewcommand{\chaptername}{Κεφάλαιο}
    \renewcommand{\appendixname}{Παράρτημα}
}
